\section{Discussion}\label{sec:discussion}

The central implication of the results is that endogenous reference formation changes both what a CPT portfolio learner optimizes and how that objective should be learned.  E1 shows that matched accounts can rank the same feasible portfolios differently solely because their wealth histories imply different references.  Once this behavioral state is allowed to evolve, the learner faces a moving objective whose gradient also contains a nonlinear distributional term.  The estimator and online analyses address these two consequences separately but within one first-order framework: the frozen-target gradient can be estimated exactly without finite-sample bias under the complete-trajectory oracle, and the resulting projected updates can be controlled through the contemporaneous residual and dynamic local regret.

The controlled experiments make this decomposition visible.  The estimator study isolates finite-budget behavior of the gradient oracle, while the online experiment shows how continued updating changes the accumulation of projected first-order error relative to a frozen policy.  The reference-mechanism experiment adds a behavioral interpretation to this separation.  Online adaptation remains useful even when reference updating is removed, so the tracking advantage is not created by asymmetry alone; at the same time, gain-dominant reference adaptation strengthens the separation and shapes how the moving target evolves.  This distinction matters because it locates the role of the behavioral mechanism more precisely: the reference dynamics structure the objective being tracked, whereas online updating determines how well the learner follows that objective.

The market and behavioral experiments suggest that the same mechanism remains informative beyond the controlled setting.  In the chronological A-share replay, DRCPT-PG has the highest mean Sharpe ratio and the lowest mean drawdown, turnover, and transaction cost across the five policy-gradient methods.  In held-out investor data, introducing a reference produces the main predictive gain, while adaptive and asymmetric updating provide further resolved improvements and enlarge the reference-conditioned disposition gap.  The heterogeneous-agent replay then shows a different consequence: adoption rates remain similar across methods, but DRCPT-PG ranks first on adoption-weighted interaction value in every calibrated investor type, indicating that the difference is driven by what happens after adoption rather than by uniformly higher acceptance.

These findings also clarify the scope of the framework.  The theoretical guarantees concern the regularized frozen-target gradient and online tracking under controlled objective variation and simulator error; they do not require downstream market performance to serve as a proxy for first-order convergence.  Conversely, the market and behavioral studies address questions the theorems do not answer: whether the reference mechanism remains useful in point-in-time equity data, whether it improves prediction of held-out investor decisions, and how reference-aware recommendations interact with heterogeneous behavioral types.  The combination is therefore most informative when read as a sequence from mechanism to optimization and then to external transfer.

Several extensions follow naturally from this structure.  On the statistical side, richer debiasing constructions could reduce the trajectory cost of the frozen-target oracle while preserving exactness.  On the online side, sharper objective-sensitivity conditions could replace the present slow-variation regime with adaptation guarantees under more abrupt market and reference shifts.  On the behavioral side, the next step is a prospective human-in-the-loop study in which reference adaptation, recommendation uptake, realized portfolio outcomes, and user-reported utility are collected under a common pre-registered protocol.  Such data would allow the dynamic-reference mechanism to be tested jointly as an optimization state and as a behavioral description of how investors respond to evolving gains and losses.
